% =============================================================
% บทที่ 5 การใช้งานสำหรับผู้ดูแลระบบ (หน้าผู้ดูแลระบบบนเว็บ)
% เนื้อหาเรียงตาม Use Case ในบทที่ 4 ของเล่มรายงาน
% ตรวจกับเว็บจริงที่รันบน localhost และซอร์สโค้ดใน main repo/frontend (27 ต.ค. 2569)
% =============================================================
\label{ch:admin-web}

\section{การจัดการสถานที่}
\label{sec:admin-places}

หน้าจัดการสถานที่แสดงดังภาพที่~\ref{fig:admin-places}

\begin{screenshot}[width=\linewidth]{images/screens/admin-places.jpg}{หน้าจัดการสถานที่}{fig:admin-places}
  \calloutline{0.48}{0.085}{0.48}{0.125}{1}
  \calloutline{0.59}{0.085}{0.59}{0.125}{2}
  \calloutline{0.71}{0.085}{0.71}{0.125}{3}
  \calloutline{0.83}{0.085}{0.83}{0.125}{4}
  \calloutline{0.925}{0.085}{0.925}{0.125}{5}
  \callout{0.80}{0.245}{6}
  \callout{0.345}{0.545}{7}
  \calloutline{0.205}{0.69}{0.205}{0.645}{8}
  \callout{0.30}{0.69}{9}
\end{screenshot}

\begin{calloutlist}
  \item ปุ่ม \btn{ซิงก์จาก Google} (หัวข้อ~\ref{sec:place-sync})
  \item ปุ่ม \btn{ตั้งค่าซิงก์อัตโนมัติ} (หัวข้อ~\ref{sec:auto-sync})
  \item ปุ่ม \btn{เพิ่มจาก Google Maps}
  \item ปุ่ม \btn{นำเข้าตามพื้นที่}
  \item ปุ่ม \btn{เพิ่มสถานที่ใหม่}
  \item ช่องค้นหาและตัวกรองตามสถานะ (ซ่อนอยู่หรือฉบับร่าง และเผยแพร่แล้ว) อำเภอ และตำบล พร้อมการเรียงตามชื่อหรือคะแนน
  \item วันที่ซิงก์ล่าสุดของสถานที่ (``ยังไม่เคยซิงก์'' ถ้ายังไม่เคยซิงก์)
  \item ปุ่ม \btn{ดู} \btn{แก้ไข} และ \btn{ซ่อนจากหน้าเว็บ}
  \item ปุ่มจุดสามจุด มีคำสั่ง \btn{ลบ}
\end{calloutlist}

\subsection{การเพิ่มสถานที่}

การเพิ่มสถานที่ทำได้ 3 วิธี ดังนี้

\begin{enumerate}
\item ปุ่ม \btn{เพิ่มสถานที่ใหม่} กรอกข้อมูลเองในแบบฟอร์มดังภาพที่~\ref{fig:admin-place-form} ได้แก่ ชื่อ หมวดหมู่ ที่อยู่ พิกัดที่เลือกจากแผนที่ สิ่งอำนวยความสะดวก แท็ก เวลาทำการ ช่วงราคา และรูปภาพ เมื่อกรอกชื่อและเลือกหมวดหมู่แล้ว คลิกปุ่ม \btn{สร้างคำอธิบายด้วย AI} ให้ระบบสรุปคำอธิบายจากรีวิว Google Maps และผลค้นหา Google แล้วคลิกปุ่ม \btn{ใช้ข้อความนี้} หลังตรวจสอบ
\item ปุ่ม \btn{เพิ่มจาก Google Maps} ค้นหาสถานที่จากชื่อแล้วเลือกผลลัพธ์ ระบบดึงข้อมูลและรูปภาพไม่เกิน 5 รูปมาสร้างเป็นฉบับร่างที่ซ่อนอยู่ให้ตรวจแก้ก่อนเผยแพร่ สถานที่ที่มีในระบบแล้วมีป้าย ``มีในระบบแล้ว''
\item ปุ่ม \btn{นำเข้าตามพื้นที่} พิมพ์ชื่อพื้นที่ เช่น ตำบล อำเภอ หรือย่าน เลือกรัศมี จำนวนรีวิวขั้นต่ำ และหมวดหมู่ (คาเฟ่ ร้านอาหาร สถานที่ท่องเที่ยว ช้อปปิ้ง) แล้วคลิกค้นหา (ภาพที่~\ref{fig:admin-area-import}) ระบบแสดงสถานที่หมวดละไม่เกิน 20 แห่ง เลือกแห่งที่ต้องการแล้วคลิกนำเข้า โดยเลือกได้ว่าจะให้ AI สร้างคำอธิบายและจัดหมวดหมู่ และจะนำเข้าเป็นฉบับร่างหรือเผยแพร่ทันที แห่งที่ AI ไม่แน่ใจหมวดหมู่จะคงเป็นฉบับร่าง คลิกปุ่ม \btn{ยกเลิกที่เหลือ} เพื่อหยุดระหว่างนำเข้า
\end{enumerate}

\appscreen{admin-place-form}{หน้าฟอร์มเพิ่มและแก้ไขข้อมูลสถานที่}

\appscreen{admin-area-import}{หน้าต่างนำเข้าสถานที่ตามพื้นที่}

\begin{note}
การนำเข้าและการซิงก์กับ Google Places มีค่าใช้จ่ายตามจำนวนคำขอ ระบบจึงแสดงค่าใช้จ่ายโดยประมาณไว้ที่ปุ่ม
\end{note}

\subsection{การซ่อนและการลบสถานที่}

\btn{ซ่อนจากหน้าเว็บ} ซ่อนสถานที่จากผู้ใช้และจากการค้นหาของแชตบอตโดยไม่ลบข้อมูล ส่วน \btn{ลบ} ในเมนูปุ่มจุดสามจุดลบสถานที่ถาวร

\begin{warning}
การลบสถานที่กู้คืนไม่ได้ และลบ QR Code ที่ผูกอยู่และรูปภาพของสถานที่ไปด้วย ถ้าต้องการเพียงไม่ให้แสดงผล ให้ใช้การซ่อนแทน
\end{warning}

\section{การซิงก์ข้อมูลสถานที่จาก Google}
\label{sec:place-sync}

การซิงก์ดึงข้อมูลล่าสุดจาก Google Places~\cite{googleplaces2026} มาปรับปรุงสถานที่ที่มีอยู่ หน้าต่างซิงก์แสดงดังภาพที่~\ref{fig:admin-place-sync}

\appscreen{admin-place-sync}{หน้าต่างซิงก์ข้อมูลสถานที่จาก Google}

\begin{steps}
  \item ที่หน้าจัดการสถานที่ คลิกปุ่ม \btn{ซิงก์จาก Google}
  \item เลือกขอบเขต ได้แก่ ทุกสถานที่ที่มาจาก Google ตามเงื่อนไข (เช่น ไม่ได้ซิงก์เกินจำนวนวันที่กำหนด หรือมีรายงานจากผู้ใช้ค้างอยู่) หรือเลือกสถานที่เอง
  \item เลือกฟิลด์ที่ต้องการปรับปรุง
  \item คลิกปุ่ม \btn{เริ่มซิงก์} หรือ \btn{ทดลองรัน} เพื่อดูผลต่างโดยไม่บันทึกลงฐานข้อมูล
        \result{ระบบแสดงความคืบหน้า คลิกปุ่ม \btn{ยกเลิกงานที่เหลือ} เพื่อหยุดระหว่างทำงานได้}
\end{steps}

\begin{note}
ชุดงานหนึ่งรองรับไม่เกิน 500 แห่ง ฟิลด์ที่ผู้ดูแลระบบแก้ไขเองถูกล็อกและไม่ถูกเขียนทับ ถ้า Google มีค่าใหม่ในฟิลด์ที่ล็อก ระบบพักค่านั้นไว้ให้เลือก \btn{ใช้ค่าจาก Google} หรือ \btn{คงค่าเดิม} เป็นรายฟิลด์
\end{note}

\section{การตั้งค่าการซิงก์อัตโนมัติ}
\label{sec:auto-sync}

ระบบซิงก์ข้อมูลจาก Google เองทุกวันได้ หน้าต่างตั้งค่าแสดงจำนวนคำขอที่ใช้ไปในเดือนนี้ ค่าใช้จ่ายโดยประมาณ จำนวนสถานที่ที่ถึงรอบ เวลารันครั้งถัดไป และผลการรันครั้งล่าสุด ดังภาพที่~\ref{fig:admin-auto-sync}

\appscreen{admin-auto-sync}{หน้าต่างตั้งค่าการซิงก์อัตโนมัติ}

\begin{steps}
  \item ที่หน้าจัดการสถานที่ คลิกปุ่ม \btn{ตั้งค่าซิงก์อัตโนมัติ}
  \item กำหนดค่า
        \begin{enumerate}
          \item เปิดหรือปิดการซิงก์ตามเวลา และชั่วโมงที่เริ่มรันของแต่ละวัน
          \item รอบการซิงก์ของสถานที่ทั่วไป (ค่าเริ่มต้น 30 วัน) และของสถานที่สำคัญ (ค่าเริ่มต้น 7 วัน) โดยเลือกได้ว่าให้นับสถานที่แนะนำ สถานที่ที่มี QR และสถานที่ที่มีรายงานค้างอยู่เป็นสถานที่สำคัญ
          \item คำขอสูงสุดต่อเดือน (ค่าเริ่มต้น 900) และจำนวนสูงสุดต่อวัน (ค่าเริ่มต้น 100 แห่ง)
          \item ซิงก์สถานที่ที่ซ่อนอยู่ด้วย และซ่อนสถานที่ที่ปิดถาวรอัตโนมัติ (ไม่ลบข้อมูล และเปิดกลับได้ในหน้าจัดการสถานที่)
        \end{enumerate}
  \item คลิกปุ่ม \btn{บันทึกการตั้งค่า}
        \result{ระบบซิงก์สถานที่ที่ถึงรอบตามเวลาที่กำหนด}
\end{steps}

ระบบข้ามการรันของวันนั้นเมื่อใช้คำขอครบเพดานของเดือน ไม่มีสถานที่ถึงรอบ หรือมีงานซิงก์อื่นทำงานอยู่ คลิกปุ่ม \btn{ซิงก์ตอนนี้} เพื่อรันตามเงื่อนไขเดียวกันโดยไม่ต้องรอเวลา

\section{การจัดการกิจกรรม}

หน้าจัดการกิจกรรมแสดงช่วงวันจัดงาน สถานะ และสถานะในดัชนีค้นหาของแชตบอต ดังภาพที่~\ref{fig:admin-events} กิจกรรมที่จบแล้วไม่อยู่ในดัชนีค้นหา

\appscreen{admin-events}{หน้าจัดการกิจกรรม}

\begin{steps}
  \item คลิกเมนู \btn{กิจกรรม} แล้วคลิกปุ่ม \btn{สร้างกิจกรรม}
  \item กรอก \field{ชื่อ} \field{ประเภท} \field{รายละเอียด} \field{ช่วงวันที่จัดงาน} \field{สถานที่จัดงาน} \field{ค่าเข้างาน} \field{ผู้จัดงาน} และ \field{กลุ่มที่เหมาะสม}
  \item คลิกบันทึก
        \result{กิจกรรมแสดงบนหน้าเว็บ โดยระบบกำหนดสถานะ เร็ว ๆ นี้ กำลังจัดอยู่ หรือจบแล้ว จากวันที่จัดงาน}
\end{steps}

\begin{tip}
วางข้อความจากโพสต์ Facebook ในช่องด้านบนของแบบฟอร์ม แล้วคลิกปุ่ม \btn{ดึงข้อมูลอัตโนมัติ} ระบบแยกข้อมูลเป็นฟิลด์และแปลงปี พ.ศ. เป็นวันที่มาตรฐานให้ (ภาพที่~\ref{fig:admin-event-extract}) ตรวจข้อมูลก่อนบันทึกทุกครั้ง ถ้าดึงข้อมูลไม่สำเร็จ ยังกรอกแบบฟอร์มเองได้ตามปกติ
\end{tip}

\appscreen{admin-event-extract}{แบบฟอร์มกิจกรรมพร้อมเครื่องมือดึงข้อมูลจากข้อความโพสต์ด้วย AI}

ถ้างานถูกยกเลิก ให้เปลี่ยนสถานะเป็น ``ยกเลิกกิจกรรม'' หน้าเว็บจะแสดงป้าย ``ยกเลิก''

\section{การจัดการฐานความรู้แชตบอต}
\label{sec:admin-knowledge}

ฐานความรู้คือเอกสารที่แชตบอตใช้อ้างอิงเมื่อตอบคำถาม หน้าจัดการแสดงดังภาพที่~\ref{fig:admin-knowledge}

\appscreen{admin-knowledge}{หน้าจัดการฐานความรู้แชตบอต}

\begin{steps}
  \item คลิกเมนู \btn{ฐานความรู้แชทบอท} แล้วคลิกปุ่ม \btn{เพิ่มฐานความรู้}
  \item กรอก \field{หัวข้อ} \field{เนื้อหา} และ \field{หมวดหมู่} กำหนดว่าให้ใช้งานอยู่หรือไม่ และปักหมุดเป็นข้อมูลสำคัญหรือไม่ แล้วบันทึก
  \item ที่หน้าแดชบอร์ด คลิกปุ่ม \btn{อัปเดตดัชนีค้นหา}
        \result{รายการแสดงสถานะว่าอยู่ในดัชนีค้นหาแล้ว และแชตบอตค้นพบข้อมูลใหม่}
\end{steps}

\begin{important}
ทุกครั้งที่เพิ่มหรือแก้ไขสถานที่ กิจกรรม หรือฐานความรู้ ต้องอัปเดตดัชนีค้นหา มิฉะนั้นแชตบอตจะยังตอบจากข้อมูลเดิม
\end{important}

\section{การตรวจรายงานข้อมูล}
\label{sec:admin-reports}

เมนู \btn{รายงานข้อมูล} แสดงคิวรายงานที่ผู้ใช้แจ้งว่าข้อมูลไม่ถูกต้อง แยกเป็นแท็บสถานที่และกิจกรรม และจัดกลุ่มตามสถานที่หรือกิจกรรมและหัวข้อ ดังภาพที่~\ref{fig:admin-reports} สถานะของรายงานแสดงดังตารางที่~\ref{tab:report-status}

\appscreen{admin-reports}{หน้ารายงานข้อมูล}

\begin{table}[H]
\caption{สถานะของรายงานข้อมูล}
\label{tab:report-status}
\centering
\begin{tabular}{|L{3.6cm}|L{9.6cm}|}
\hline
\textbf{สถานะ} & \textbf{ความหมาย} \\ \hline
รอตรวจสอบ & รายงานใหม่ที่ยังไม่มีผู้ดูแลระบบดำเนินการ \\ \hline
จัดการแล้ว & ผู้ดูแลระบบแก้ไขข้อมูลหรือยืนยันว่าข้อมูลถูกต้องแล้ว \\ \hline
Google อัปเดตแล้ว & ปิดรายงานด้วยข้อมูลล่าสุดจากการซิงก์กับ Google \\ \hline
ปฏิเสธ & รายงานไม่ถูกต้อง \\ \hline
\end{tabular}
\end{table}

\begin{steps}
  \item คลิกเมนู \btn{รายงานข้อมูล} แล้วเลือกรายงานที่สถานะรอตรวจสอบ
  \item เทียบข้อมูลในระบบกับหมายเหตุของผู้แจ้ง แล้วเลือกดำเนินการ
        \begin{enumerate}
          \item ปุ่ม \btn{ซิงก์จาก Google} ดึงข้อมูลล่าสุดของสถานที่นั้นมาเทียบ ถ้า Google มีค่าใหม่ในฟิลด์ที่ล็อก ให้เลือก \btn{ใช้ค่าจาก Google} หรือ \btn{คงค่าเดิม}
          \item ปุ่ม \btn{เปิดหน้าแก้ไข} แก้ไขข้อมูลเอง แล้วปิดรายงานด้วยผลการแก้ไข
          \item ปุ่ม \btn{ข้อมูลถูกต้องแล้ว} ปิดรายงานโดยไม่แก้ไข หรือ \btn{ปฏิเสธ} เมื่อรายงานไม่ถูกต้อง
        \end{enumerate}
        \result{รายงานเปลี่ยนสถานะตามตารางที่~\ref{tab:report-status}}
\end{steps}

\section{การตรวจคำขอกิจกรรม}
\label{sec:admin-event-requests}

เมนู \btn{คำขอกิจกรรม} แสดงกิจกรรมที่สมาชิกส่งเข้ามา กรองตามสถานะรออนุมัติ อนุมัติแล้ว และไม่อนุมัติ ดังภาพที่~\ref{fig:admin-event-requests}

\appscreen{admin-event-requests}{หน้าคำขอกิจกรรม}

\begin{steps}
  \item คลิกเมนู \btn{คำขอกิจกรรม} แล้วเปิดดูรายละเอียดและรูปภาพของคำขอ
  \item คลิกปุ่ม \btn{อนุมัติ} หรือ \btn{ไม่อนุมัติ} พร้อมระบุเหตุผลไม่เกิน 500 ตัวอักษร
        \result{ถ้าอนุมัติ ระบบสร้างกิจกรรมจากคำขอและแสดงบนหน้าเว็บทันที ถ้าไม่อนุมัติ ผู้ขอเห็นเหตุผลที่ระบุ}
\end{steps}

ถ้าคำขอถูกตรวจไปแล้วโดยผู้ดูแลระบบคนอื่น ระบบแจ้งว่าดำเนินการซ้ำไม่ได้

\section{การนำเข้าสถานที่จาก Google ที่แชตบอตแนะนำ}

เมนู \btn{สถานที่จาก Google} แสดงสถานที่ที่แชตบอตค้นเพิ่มจาก Google เพราะข้อมูลในระบบมีน้อย เรียงตามจำนวนครั้งที่แชตบอตแนะนำ พร้อมพื้นที่ที่ผู้ใช้ถามถึง ประเภท คะแนน Google ที่อยู่ และลิงก์ \btn{ดูใน Google Maps} ด้านบนแสดงจำนวนครั้งที่แชตบอตค้น Google ของวันนี้และเดือนนี้พร้อมค่าใช้จ่ายโดยประมาณ ดังภาพที่~\ref{fig:admin-place-candidates}

\appscreen{admin-place-candidates}{หน้าสถานที่จาก Google}

\begin{steps}
  \item คลิกเมนู \btn{สถานที่จาก Google}
  \item คลิกปุ่ม \btn{นำเข้า} ที่สถานที่ที่ต้องการ
        \result{ระบบดึงรายละเอียดและรูปภาพจาก Google ให้ AI ร่างคำอธิบาย แล้วตรวจคุณภาพ}
\end{steps}

การตรวจคุณภาพดูว่ามีคำอธิบาย พิกัด รูปภาพ เวลาเปิดปิด (ยกเว้นวัด) และอำเภอ ยังเปิดกิจการ และอยู่ในจังหวัดขอนแก่น ถ้าผ่านทุกข้อ ระบบเผยแพร่และอัปเดตดัชนีค้นหาให้ ถ้าไม่ผ่าน สถานที่จะเป็นฉบับร่างที่ซ่อนอยู่พร้อมรายการปัญหาให้แก้ก่อนเผยแพร่ สถานที่ที่มีในระบบแล้วจะเชื่อมโยงกับสถานที่เดิมโดยไม่นำเข้าซ้ำ รายการที่ไม่ต้องการให้คลิกปุ่ม \btn{ไม่สนใจ} และนำกลับเข้าคิวได้จากแท็บไม่สนใจ

\section{การจัดการ QR Code}

เมนู \btn{QR Code} ใช้สร้าง QR Code ผูกกับสถานที่ ดังภาพที่~\ref{fig:admin-qr} การ์ดแต่ละรายการแสดงจำนวนครั้งที่สแกนและพอยท์ที่แจกไป

\appscreen{admin-qr}{หน้าจัดการ QR Code}

\begin{steps}
  \item คลิกเมนู \btn{QR Code} แล้วคลิกปุ่ม \btn{สร้าง QR ใหม่}
  \item เลือก \field{สถานที่} ที่ยังไม่มี QR กำหนด \field{จำนวนพอยท์} \field{รัศมีที่สแกนได้} (20 ถึง 5,000 เมตร ดูขอบเขตบนแผนที่ได้) และ \field{วันหมดอายุ} (ถ้ามี) แล้วบันทึก
  \item คลิกปุ่ม \btn{ดู QR Code} เพื่อดาวน์โหลดหรือพิมพ์ป้ายไปติดที่สถานที่
\end{steps}

หนึ่งสถานที่มี QR ได้เพียงอันเดียว ใช้ปุ่ม \btn{แก้ไข} และปุ่มเปิดหรือปิดใช้งานเพื่อจัดการ QR ที่มีอยู่

\section{การจัดการของรางวัล}

เมนู \btn{ของรางวัล} ใช้เพิ่มและแก้ไขของรางวัล ดังภาพที่~\ref{fig:admin-rewards}

\appscreen{admin-rewards}{หน้าจัดการของรางวัล}

\begin{steps}
  \item คลิกเมนู \btn{ของรางวัล} แล้วคลิกปุ่ม \btn{เพิ่มของรางวัล}
  \item กรอก \field{ชื่อ} \field{จำนวนพอยท์ที่ใช้แลก} และ \field{จำนวนสต็อก} (เว้นว่างถ้าไม่จำกัดจำนวน) แนบรูปภาพถ้ามี แล้วบันทึก
\end{steps}

\section{การแลกของรางวัลที่เคาน์เตอร์}
\label{sec:admin-redeem}

เมื่อสมาชิกมาแลกของรางวัล เจ้าหน้าที่ใช้เมนู \btn{แลกที่เคาน์เตอร์} ดังภาพที่~\ref{fig:admin-redeem}

\appscreen{admin-redeem}{หน้าแลกของรางวัลที่เคาน์เตอร์}

\begin{steps}
  \item คลิกเมนู \btn{แลกที่เคาน์เตอร์}
  \item ค้นหาผู้ใช้ด้วยชื่อ อีเมล หรือเบอร์โทรศัพท์
  \item เลือกของรางวัล แล้วคลิกยืนยันเมื่อพร้อมมอบของ
        \result{ระบบหักพอยท์และลดสต็อกในขั้นตอนเดียว พร้อมบันทึกชื่อเจ้าหน้าที่ผู้ทำรายการ}
\end{steps}

ถ้าพอยท์ไม่พอ ของรางวัลหมด หรือบัญชีถูกระงับ ระบบแจ้งสาเหตุและไม่บันทึกการแลก ตารางประวัติการแลกด้านล่างกรองได้ตามสถานะแลกแล้วและยกเลิกแล้ว ถ้าบันทึกผิด ให้คลิกปุ่ม \btn{ยกเลิกรายการ} พร้อมระบุเหตุผล ระบบคืนพอยท์ให้ผู้ใช้และคืนของรางวัลเข้าสต็อก

\section{การจัดการบัญชีผู้ใช้}
\label{sec:admin-users}

เมนู \btn{ผู้ใช้} ใช้ค้นหาด้วยชื่อ อีเมล หรือเบอร์โทรศัพท์ กรองตามสถานะ (ใช้งานอยู่ ถูกระงับ ลบแล้ว) บทบาท และการยืนยันอีเมล และเรียงตามวันที่สมัคร ชื่อ พอยท์ หรือการใช้งานล่าสุด ดังภาพที่~\ref{fig:admin-users}

\appscreen{admin-users}{หน้าจัดการผู้ใช้}

คลิกผู้ใช้เพื่อดูรายละเอียดบัญชี ประวัติการสแกน ประวัติการแลก ประวัติการปรับพอยท์ และประวัติการดำเนินการของผู้ดูแลระบบ ดังภาพที่~\ref{fig:admin-user-detail} ปุ่มดำเนินการมีดังนี้

\appscreen{admin-user-detail}{หน้ารายละเอียดและการจัดการบัญชีผู้ใช้}

\begin{enumerate}
\item ปุ่ม \btn{ระงับบัญชี} หรือ \btn{เปิดใช้งานบัญชี}
\item ปุ่ม \btn{ลบบัญชี} หรือ \btn{กู้คืนบัญชี}
\item ปุ่ม \btn{ปรับพอยท์} ระบุจำนวนที่เพิ่มหรือหักและเหตุผล
\item ปุ่ม \btn{ออกจากระบบทุกอุปกรณ์}
\item ปุ่ม \btn{ส่งอีเมลยืนยันใหม่} สำหรับบัญชีที่ยังไม่ยืนยันอีเมล
\end{enumerate}

\begin{note}
ทุกการดำเนินการถูกบันทึกพร้อมชื่อผู้ดำเนินการและเวลา และระบบไม่อนุญาตให้ระงับหรือลบบัญชีผู้ดูแลระบบคนสุดท้าย
\end{note}

\section{การดูทริปของผู้ใช้}

เมนู \btn{ทริป} แสดงแผนทริปที่ผู้ใช้สร้าง กรองได้ตามแผนของผู้ที่เข้าสู่ระบบ ผู้ที่ไม่ได้เข้าสู่ระบบ และแผนที่มีสถานที่ปิดหรือถูกลบ ดังภาพที่~\ref{fig:admin-trips} ใช้ติดตามการใช้งานและตรวจคุณภาพของแผน

\appscreen{admin-trips}{หน้าดูทริปของผู้ใช้}

\section{การดูผลความพึงพอใจ}

เมนู \btn{ความพึงพอใจ} แสดงผลแบบสอบถาม ได้แก่ จำนวนผู้ตอบ ค่าเฉลี่ยรวมและส่วนเบี่ยงเบนมาตรฐาน ระดับความพึงพอใจโดยรวม จำนวนคำตอบที่มีข้อเสนอแนะ สรุปรายด้านและรายข้อ และคำตอบรายคน ดังภาพที่~\ref{fig:admin-survey} ข้อที่ผู้ตอบเลือก ``ยังไม่เคยใช้'' ไม่ถูกนำมาคำนวณ

\appscreen{admin-survey}{หน้าผลความพึงพอใจ}

คลิกปุ่ม \btn{ดาวน์โหลด CSV} เพื่อส่งออกข้อมูล และลบคำตอบที่เป็นข้อมูลทดสอบหรือไม่เหมาะสมได้ คำตอบที่ลบจะไม่ถูกนับในสถิติ

\section{การเพิ่มและอัปเดตฐานข้อมูล}
\label{sec:data-update}

ผู้ดูแลระบบมีหน้าที่เพิ่มข้อมูลใหม่และอัปเดตข้อมูลสถานที่ กิจกรรม และฐานความรู้ให้ตรงกับความเป็นจริงอยู่เสมอ ข้อมูลเข้าสู่ระบบจาก 3 ทาง ได้แก่ ผู้ดูแลระบบกรอกเอง Google Maps และคิวตรวจสอบที่เกิดจากการใช้งานของนักท่องเที่ยว โดยข้อมูลทุกทางต้องผ่านการตรวจสอบของผู้ดูแลระบบก่อนแสดงบนหน้าเว็บ ภาพที่~\ref{fig:upd-cycle} แสดงวงจรการเพิ่มและอัปเดตฐานข้อมูล

\begin{figure}[H]
\centering
\begin{tikzpicture}[
  box/.style={draw=black, fill=white, rounded corners=3pt, thick,
              text width=4.4cm, align=center, minimum height=1.9cm, inner sep=6pt,
              font=\normalfont\small},
  arr/.style={-{Stealth[length=3mm]}, thick, draw=black!60},
  lbl/.style={font=\normalfont\footnotesize, text=black!75, inner sep=2bp}]
\node[box] (in)  at (0,4.4)  {เพิ่มข้อมูลใหม่\\กรอกเอง หรือนำเข้าจาก Google Maps};
\node[box] (rev) at (0,1.6)  {ตรวจทานก่อนเผยแพร่\\เก็บเป็นฉบับร่างที่ซ่อนอยู่จนตรวจสอบแล้ว};
\node[box] (pub) at (8.2,1.6)  {เผยแพร่และอัปเดตดัชนีค้นหา\\ข้อมูลแสดงบนหน้าเว็บและแชตบอต};
\node[box] (que) at (8.2,-1.8) {ตรวจคิว\\รายงานข้อมูล คำขอกิจกรรม และสถานที่จาก Google};
\node[box] (upd) at (0,-1.8) {อัปเดตข้อมูล\\ซิงก์กับ Google และแก้ไขตามคิว};
\draw[arr] (in) -- (rev);
\draw[arr] (rev) -- node[lbl, above]{เผยแพร่} (pub);
\draw[arr] (pub) -- (que);
\draw[arr] (que) -- (upd);
\draw[arr] (upd) -- (rev);
\end{tikzpicture}
\caption{วงจรการเพิ่มและอัปเดตฐานข้อมูล}
\label{fig:upd-cycle}
\end{figure}

\subsection{การเพิ่มข้อมูลใหม่}

ผู้ดูแลระบบเพิ่มสถานที่ได้ 3 วิธี คือ กรอกเอง นำเข้าทีละแห่งจาก Google Maps หรือนำเข้าครั้งละหลายแห่งตามพื้นที่ (หัวข้อ~\ref{sec:admin-places}) เพิ่มกิจกรรมจากเมนูกิจกรรม หรืออนุมัติคำขอกิจกรรมที่สมาชิกส่งเข้ามา (หัวข้อ~\ref{sec:admin-event-requests}) และเพิ่มเอกสารให้แชตบอตจากเมนูฐานความรู้แชทบอท (หัวข้อ~\ref{sec:admin-knowledge})

อีกแหล่งหนึ่งคือคิว ``สถานที่จาก Google'' เมื่อนักท่องเที่ยวถามแชตบอตถึงพื้นที่ที่ระบบมีข้อมูลน้อย แชตบอตจะค้นเพิ่มจาก Google และแสดงการ์ดที่ติดป้าย ``แหล่งจาก Google'' ดังภาพที่~\ref{fig:upd-chatbot-google} ระบบเก็บสถานที่เหล่านั้นไว้ในคิวเรียงตามจำนวนครั้งที่ผู้ใช้ถามถึง ดังภาพที่~\ref{fig:upd-candidates} ผู้ดูแลระบบจึงเลือกนำเข้าสถานที่ในพื้นที่ที่ผู้ใช้สนใจจริงได้ก่อน

\appscreen[0.42]{upd-chatbot-google}{การ์ดสถานที่ที่แชตบอตค้นเพิ่มจาก Google}

\appscreen{upd-candidates}{คิวสถานที่จาก Google ที่แชตบอตแนะนำให้ผู้ใช้}

สถานที่ที่เพิ่งนำเข้าอยู่ในสถานะฉบับร่างที่ซ่อนจากผู้ใช้ ผู้ดูแลระบบตรวจทานหมวดหมู่ คำอธิบาย รูปภาพ และเวลาเปิดปิดก่อนเผยแพร่ เมื่อเผยแพร่แล้ว สถานที่จะปรากฏในหน้ารายการ ผลการค้นหา แชตบอต และแผนการเดินทางทันที

\subsection{การอัปเดตข้อมูลสถานที่}

ข้อมูลสถานที่ที่มาจาก Google Maps อัปเดตด้วยการซิงก์ ผู้ดูแลระบบสั่งซิงก์เองได้ (หัวข้อ~\ref{sec:place-sync}) หรือตั้งให้ระบบซิงก์อัตโนมัติทุกวัน (หัวข้อ~\ref{sec:auto-sync}) ดังภาพที่~\ref{fig:upd-auto-sync} สถานที่สำคัญ เช่น สถานที่แนะนำ สถานที่ที่มี QR Code และสถานที่ที่มีรายงานค้างอยู่ จะถูกซิงก์ถี่กว่าสถานที่ทั่วไป และระบบจำกัดจำนวนคำขอต่อวันและต่อเดือนเพื่อควบคุมค่าใช้จ่าย

\appscreen[0.6]{upd-auto-sync}{การกำหนดรอบการซิงก์และเพดานค่าใช้จ่ายของการซิงก์อัตโนมัติ}

ข้อมูลที่ผู้ดูแลระบบแก้ไขเองถูกล็อกไว้ การซิงก์จึงไม่เขียนทับ ถ้า Google มีค่าใหม่ในฟิลด์นั้น ระบบแสดงค่าใหม่ไว้ให้ผู้ดูแลระบบเลือกว่าจะใช้หรือคงค่าเดิม นอกจากการซิงก์ ผู้ดูแลระบบอัปเดตข้อมูลจากคิวรายงานข้อมูลที่สมาชิกแจ้งเข้ามาตามหัวข้อ~\ref{sec:admin-reports}

ผลของการอัปเดตตรวจได้ 2 จุด หน้ารายละเอียดสถานที่แสดงวันที่อัปเดตล่าสุดไว้ท้ายหน้า ดังภาพที่~\ref{fig:upd-place-updated} และหน้าจัดการสถานที่แสดงวันที่ซิงก์ล่าสุดบนการ์ดของแต่ละสถานที่ ดังภาพที่~\ref{fig:upd-place-sync-date} สถานที่ที่ปิดถาวรจะแสดงป้าย ``ปิดถาวรแล้ว'' และซ่อนจากหน้าเว็บได้อัตโนมัติถ้าเปิดตัวเลือกนี้ไว้ในการตั้งค่าการซิงก์อัตโนมัติ

\appscreen{upd-place-updated}{วันที่อัปเดตข้อมูลล่าสุดในหน้ารายละเอียดสถานที่}

\appscreen{upd-place-sync-date}{วันที่ซิงก์ล่าสุดบนการ์ดสถานที่ในหน้าผู้ดูแลระบบ}

\subsection{การอัปเดตสถานะกิจกรรม}

ระบบเปลี่ยนสถานะของกิจกรรมจากวันที่จัดงานให้อัตโนมัติ ได้แก่ เร็ว ๆ นี้ กำลังจัดอยู่ และจบแล้ว ผู้ดูแลระบบจึงไม่ต้องแก้สถานะเอง กิจกรรมที่จบแล้วจะไม่อยู่ในดัชนีค้นหาของแชตบอต ดังภาพที่~\ref{fig:upd-event-status} ผู้ดูแลระบบต้องแก้สถานะเองเฉพาะกรณีงานถูกยกเลิก โดยเปลี่ยนเป็น ``ยกเลิกกิจกรรม''

\appscreen{upd-event-status}{สถานะของกิจกรรมที่ระบบเปลี่ยนตามวันที่จัดงาน}

\subsection{การอัปเดตดัชนีค้นหาของแชตบอต}

แชตบอตค้นคำตอบจากดัชนีค้นหาของข้อมูลสถานที่ กิจกรรม และฐานความรู้ หลังเพิ่มหรือแก้ไขข้อมูล แดชบอร์ดแสดงจำนวนรายการที่ยังไม่เข้าดัชนีค้นหา ผู้ดูแลระบบคลิกปุ่ม \btn{อัปเดตดัชนีค้นหา} เพื่อให้แชตบอตค้นพบข้อมูลล่าสุด ดังภาพที่~\ref{fig:upd-index-card} ส่วนสถานที่ที่นำเข้าจากคิวสถานที่จาก Google และผ่านการตรวจคุณภาพ ระบบอัปเดตดัชนีค้นหาให้ทันที

\appscreen[0.6]{upd-index-card}{การ์ดดัชนีค้นหาแชตบอตบนหน้าแดชบอร์ด}
